\section{State-dependent noncommuting instruments}\label{sec:quantum}
Let $D\ge2$, $d=D^2-1$, and choose a real Hilbert--Schmidt orthonormal basis $F_1,\ldots,F_d$ of traceless Hermitian matrices. The prediction coordinate is $x_j=\operatorname{tr}(\rho F_j)$; its Euclidean distance is the Hilbert--Schmidt distance of density operators. This finite coordinate system will be verified by the raw instrument and audit. It is not memory available to the classical observer.

Let $\mu$ be normalized Lebesgue measure on the unit Hilbert--Schmidt ball of traceless Hermitian matrices $H$, and fix $0<\eta<1$, $0<\lambda<1$ and a known unitary $U$. Define
\[
 A_H=I+\eta H,\quad K_H=A_H^{1/2},\quad
 T(\sigma)=\lambda U\sigma U^*+(1-\lambda)\operatorname{tr}(\sigma)I/D.
\]
One raw call is the completely positive instrument
\begin{equation}\label{eq:quantum-raw}
 \mathcal I_H(\rho)=T(K_H\rho K_H)\,\mu(dH).
\end{equation}
Since $\int H\,d\mu=0$, it is normalized. Its report density and prediction update are
\[
 \ell(H,\rho)=1+\eta\operatorname{tr}(H\rho),\qquad
 F(H,\rho)=T\left(\frac{K_H\rho K_H}{\operatorname{tr}(A_H\rho)}\right).
\]
The law depends on the entering state, and different effects $A_H$ need not commute. Reports range over a Euclidean effect body, not a Haar group orbit. The physical channel and measurement together are one declared instrument call; any implementation that resolves its internal suboperations must charge their physical time separately.

Prepare a known $\rho_0\in K=T(\mathcal D_D)$, where $\mathcal D_D$ is the density-operator body. Then every prediction has eigenvalues between $a=(1-\lambda)/D$ and $b=a+\lambda$. The bounded executable audit chooses $I$ uniformly in $\{1,\ldots,d\}$ after commitment, measures the two effects $(I_D\pm F_I)/2$, and returns $Z=dS e_I$. Positivity follows from $\|F_i\|_{\rm op}\le\|F_i\|_{\rm HS}=1$. Directly,
\begin{equation}\label{eq:audit-normalization}
 \E[S e_I\mid\rho]=x/d,\qquad \E[Z\mid\rho]=x,
 \qquad \E\norm Z^2=d^2.
\end{equation}
Only one basis test is performed, not a joint measurement of incompatible effects. All losses can be rescaled by a fixed dimension-dependent constant to lie in $[0,1]$ for a comparison theorem.

\begin{lemma}[Positive congruence coordinates]\label{lem:congruence}
For positive definite $q$, the map
\[
 A\longmapsto \frac{A^{1/2}qA^{1/2}}{\operatorname{tr}(Aq)}
\]
is a smooth bijection from positive definite matrices of trace $D$ to positive definite density operators. Its inverse at $\sigma$ is
\begin{equation}\label{eq:congruence-inverse}
 B=q^{-1/2}(q^{1/2}\sigma q^{1/2})^{1/2}q^{-1/2},
 \qquad A=DB^2/\operatorname{tr}(B^2).
\end{equation}
On compact positive definite input domains all derivatives required for a change of variables are uniformly bounded, and the absolute Jacobian has a positive lower bound.
\end{lemma}
\begin{proof}
$B$ is the unique positive definite solution of $BqB=\sigma$: conjugating by $q^{1/2}$ reduces uniqueness to the positive square root. Formula~\eqref{eq:congruence-inverse} normalizes its otherwise immaterial scale. Conversely any original $A$ is recovered by that normalization. Positive square roots and inverses are smooth on the positive definite cone. Differentiating the inverse identities on the trace hyperplanes shows that the derivatives are inverse linear maps. Compactness supplies the uniform bounds and nonvanishing Jacobian. This elementary inverse is also used in the earlier positive-instrument realization; it is not claimed anew as a matrix identity.
\end{proof}

\begin{theorem}[Weak effects followed by a positivity-improving channel]\label{thm:quantum}
The raw instrument~\eqref{eq:quantum-raw}, with the declared audit, verifies Theorem~\ref{thm:main} with acquired dimension $d=D^2-1$ and certificate (P), for every $0<\lambda,\eta<1$. Constants are uniform on compact subintervals of those parameters and fixed finite $D$. The observer is classical and has no access to the physical system except the reports and the final post-commitment audit.
\end{theorem}
\begin{proof}
The unnormalized instrument is linear in $\rho$, so all report/test continuations have the required finite affine closure. The audit coordinates separate density operators; their preparations provide a continuous realization. Positivity gives $1-\eta\le\ell\le1+\eta$, a common report support and
$\|\ell(\cdot,\rho)-\ell(\cdot,\sigma)\|_{L^1}\le\eta\|\rho-\sigma\|_{\rm HS}$. Neither an eigenbasis chosen by the unknown state nor a preexisting information metric is assumed.

Use the cone metric $h(\rho,\sigma)=\log(M/m)$, where $m\sigma\preceq\rho\preceq M\sigma$ are the extremal order bounds. Positive congruence and normalization preserve this metric. It has bounded diameter $H_0=2\log(b/a)$ on $K$ and is uniformly equivalent there to the Hilbert--Schmidt distance.

Here is a direct contraction estimate for $T$. For normalized positive $\rho,\sigma$, $m\le1\le M$. Put $c=(1-\lambda)/(Db)\in(0,1)$. Since $T(\sigma)\preceq bI_D$,
\[
 [m+c(1-m)]T(\sigma)\preceq T(\rho)
 \preceq[M-c(M-1)]T(\sigma).
\]
For example the lower follows by expanding $T(\rho)-mT(\sigma)$ and retaining its $(1-\lambda)(1-m)I_D/D$ part. The upper is the same argument applied to $M\sigma-\rho$. When $h(\rho,\sigma)\le H_0$, both $m,M$ lie in $[e^{-H_0},e^{H_0}]$. The derivative of $\log((1-c)e^v+c)$ is bounded on this interval by
\[
 \vartheta=\frac{(1-c)e^{H_0}}{(1-c)e^{H_0}+c}<1.
\]
Integrating between $\log m$ and $\log M$ proves $h(T\rho,T\sigma)\le\vartheta h(\rho,\sigma)$. Applied after the metric-isometric normalized congruence, this establishes (P). It does not require $3\lambda<1$.

For fixed $\rho\in K$, Lemma~\ref{lem:congruence} applies to $A_H$ on its compact effect ball, uniformly because $\rho\succeq aI_D$ and $A_H\succeq(1-\eta)I_D$. The linear parametrization by $H$ has rank $d$. The final channel $T$ is an invertible affine map on the trace-one hyperplane, with linear part $\lambda$ times a unitary isometry. The posterior map therefore has a uniformly nonzero Jacobian. Its actual report density is bounded by $1+\eta$, so its image measure has bounded $d$-dimensional density. This proves the acquired small-ball estimate on every actual incoming row. The effect-ball boundary has zero reference mass, and no conditional renormalization is made. Global metric covers follow from Euclidean covers of the compact convex body $K$. All requirements of Theorem~\ref{thm:main} have now been verified from~\eqref{eq:quantum-raw}.
\end{proof}

\begin{corollary}[Noncontracting phases between bounded-gap improvements]\label{cor:quantum-block}
In~\eqref{eq:quantum-raw}, allow time-dependent known unitaries and channels that are either unitary ($\lambda_t=1$) or have $\lambda_t\in[\lambda_-,\lambda_+]\subset(0,1)$. Suppose every block of $L$ consecutive phases contains an improving channel, and prepare an interior state with a fixed eigenvalue lower bound. With a common effect law and $\eta$ bounded away from one, the same matching risk laws hold uniformly in the horizon. Individual normalized updates need not contract.
\end{corollary}
\begin{proof}
A normalized effect changes a lower eigenvalue bound by at worst the factor $(1-\eta)/(1+\eta)$. An improving channel restores it to at least $(1-\lambda_+)/D$. The bounded-gap condition thus gives a common positive eigenvalue floor for all declared domains, including the initial incomplete block. Choose convex domains with these phasewise floors. Congruences and unitary phases are metric isometries; each improving phase contracts by a common constant less than one by the preceding proof. Products therefore satisfy (P). The inverse-map, report-continuity and cover constants stay uniform on the resulting compact interior domains, including unitary phases. This proves the assertion without replacing the block by an uncharged reset.
\end{proof}
